% Generated draft: requires mathematical and source review.
\documentclass[11pt]{article}
\usepackage[margin=25mm]{geometry}
\usepackage{fontspec}
\setmainfont{Latin Modern Roman}
\usepackage{amsmath,amssymb,mathtools,unicode-math}
\setmathfont{Latin Modern Math}
\usepackage{enumitem,needspace,xurl,hyperref}
\setlength{\parindent}{0pt}
\setlength{\parskip}{5pt}
\setlength{\emergencystretch}{3em}
\newcommand{\N}{\mathbb{N}}
\newcommand{\Z}{\mathbb{Z}}
\newcommand{\Q}{\mathbb{Q}}
\newcommand{\R}{\mathbb{R}}
\newcommand{\C}{\mathbb{C}}
\newcommand{\F}{\mathbb{F}}
\newcommand{\A}{\mathbb{A}}
\newcommand{\PP}{\mathbb{P}}
\newcommand{\E}{\mathbb{E}}
\newcommand{\eps}{\varepsilon}
\newcommand{\dd}{\,\mathrm{d}}
\begin{document}
\section*{UnsolvedMath ID: 30003516}
\textbf{Source problem code:} catalogue entry
\subsection{Definitions and mathematical statement}
Let \(L\) denote the number of loci in the multilocus ancestry and diffusion models referred to in the supplied source. Thus \(L\) is a positive integer. Write \(m(\cdot)\) for the function described in the supplied background as the dual normalizer. The record does not define this function, its domain, or its normalization, and does not specify the ancestry processes, diffusion generators, or duality relation.

Subject to these unresolved model specifications, the source poses the following research questions:
\begin{enumerate}
\item Incorporate recombination and selection into the multilocus ancestry and diffusion models under consideration.
\item Determine \(m(\cdot)\) in closed form beyond the one-locus reversible case mentioned in the source.
\item Extend the strong-recombination diffusion-limit arguments referred to for \(L=2\) to integer numbers of loci \(L>2\), together with a posterior dual.
\end{enumerate}
The record does not specify the meaning of a posterior dual, the conditioning information involved, the strong-recombination scaling, or the asserted limiting process and mode of convergence. It also does not identify the parameter regimes or hypotheses under which these extensions are sought. Consequently, these are source-level research questions requiring further specification, rather than a fully determined theorem or conjecture. In particular, the record does not establish whether the requested extensions must hold for every \(L>2\), every selection and recombination mechanism, or only particular model classes.
\subsection{Short English statement}
The problem asks how to include recombination and selection in the referenced multilocus ancestry and diffusion models, obtain a closed formula for the dual normalizer \(m\) beyond the one-locus reversible case, and extend the two-locus strong-recombination limit to more loci with a posterior dual. The supplied record lacks the model definitions and limiting assumptions needed for a precise mathematical formulation.
\subsection{Sources}
\begin{itemize}
\item Supplied local UnsolvedMath registry record: \texttt{data/unsolvedmath/problems.json}, permanent ID \(30003516\), source problem code \texttt{catalogue entry}, title \emph{Recombination, Selection, and Duality in Multilocus Ancestry Models}; collection: Oberwolfach Reports Open Problems. No individual proposer is supplied; the proposed year is recorded as 2017.
\item Supplied report citation: \emph{Reaction Networks and Population Dynamics (2018)}, Oberwolfach Reports, pp.\ 1747--1804, DOI \texttt{10.4171/owr/2017/28}. Supplied source URL and literature-source link, labeled ``original OWR report'': \url{https://doi.org/10.4171/owr/2017/28}. No problem-specific page or contributing author is supplied.
\item Literature assessment dated exactly \texttt{2026-08-21}: the supplied curation labels the problem open and states that the closed-form dual normalizer and the extension to more than two loci with posterior dual strong recombination remain unresolved; it further states that related multilocus inference work does not settle these questions. This is supplied dated curation, not an independently verified assessment. The background expressly requests independent verification of time-sensitive status and cited claims.
\item The supplied registry labels the statement \texttt{corrected\_verified} and describes its editing as preserving mathematical content while repairing notation, typography, or omitted definitions. This is a supplied editorial qualification, not independent mathematical verification of this draft.
\end{itemize}
\end{document}
